\chapter{Conclusion}
\label{ch:conclusion}

\section{Summary of Contributions}
\label{sec:summary}

This thesis presented a comprehensive study on \kvp{} extraction from business
documents using LayoutLMv3 with a learned span-level linker. The main
contributions are:

\begin{enumerate}
    \item \textbf{Dataset Analysis}: Characterised the KVP10k dataset through
    clustering of 9{,}124 unique coordinate-bearing training pages, yielding
    two weakly separated annotation-geometry regimes, and through a separate
    exploratory spatial-link audit (Chapter~\ref{ch:dataset}).

    \item \textbf{Generative Baseline}: Re-implemented the KVP10k Mistral-7B
    approach and identified unsupported generation and weak spatial grounding
    (Chapter~\ref{ch:exresults}).

    \item \textbf{Discriminative Architecture}: Implemented a LayoutLMv3-based
    entity classifier achieving F1~$= 0.827$. Grounded token classification
    prevents generation of text absent from the OCR input
    (Chapter~\ref{ch:methodology}).

    \item \textbf{Linking Granularity Analysis}: Demonstrated that token-level
    linking fails due to extreme class imbalance ($\sim$450:1), yielding link
    F1~$= 0.008$. Designed and implemented span-level linking (V2) to address
    this mismatch (Chapter~\ref{ch:methodology}).

    \item \textbf{Pipeline Debugging and V4 Results}: Diagnosed three
    compounding data-pipeline bugs (missing teacher forcing, over-strict
    bounding-box filter, cross-product link label explosion) that masked
    span-level linking performance. The corrected V4 model achieves regular-KVP
    macro F1~$=0.334$ (text) / $0.309$ (text+location), below the corresponding
    published scores ($0.659/0.611$); diagnostic analyses identify relation
    modeling within the current biaffine-linker pipeline as the main remaining
    bottleneck (Chapters~\ref{ch:exresults}--\ref{ch:discussion}).

    \item \textbf{Implementation}: Released a modular, reproducible codebase
    with CONFIG-controlled experiments.
\end{enumerate}

\section{Research Questions Revisited}
\label{sec:rq_revisited}

\subsection{RQ1: Layout Characteristics}
Page-level clustering identifies two weakly separated annotation-geometry
regimes. Cluster~0 has a higher box count and broader spatial spread, whereas
Cluster~1 has fewer, larger boxes in a more compact region. Two principal
components explain 52.8\% of standardized-feature variance, and the silhouette
score is 0.222. These groups are not document genres or OCR-density classes.
In a separate 200-row exploratory audit, linked \kvps{} have mean normalized
center distance 0.120, suggesting that spatial features may be informative for
linking without establishing a causal requirement.

\subsection{RQ2: Discriminative vs.\ Generative}
The discriminative LayoutLMv3 model achieves entity F1~$=0.827$ without
generating text outside the OCR input. For regular pair extraction, V4 reaches
combined macro F1 $=0.309$, compared with $0.521$ for reconstructed Mistral;
the latter nevertheless produces substantial unsupported text.
These are complementary trade-offs rather than evidence that one formulation
dominates both entity detection and linking.

\subsection{RQ3: Linking Granularity}
Token-level linking (V1) fails decisively (F1~$= 0.008$) due to the
$\sim$450:1 class imbalance inherent in token-to-token scoring. Span-level
linking (V2) reduces this to $\sim$5:1 by grouping contiguous tokens before
scoring. However, three data-pipeline bugs initially masked this benefit,
producing link F1 of only 0.007. After systematic diagnosis and correction,
the V4 model achieves regular-KVP macro F1~$=0.334$ (text) / $0.309$
(text+location).
Span-level linking, paired with correct supervision, is far more effective
than token-level linking, but still does not reach the
published regular-KVP scores ($0.659/0.611$); an internal oracle-entity
diagnostic shows that, within the current pipeline, relation scoring rather
than entity detection accounts for most of the observed gap.

\subsection{RQ4: Success and Failure Factors}
Performance varies across the exploratory annotation-geometry strata, but the
direction differs by model: Stage~3 regular-pair F1 is higher in Cluster~0,
whereas V4 regular-pair F1 is higher in Cluster~1. This association does not
establish a causal document-type effect. Diagnostics identify relation modeling within the current
biaffine-linker pipeline as the main remaining bottleneck: supplying
ground-truth entities raises pooled diagnostic link F1 by only
$\sim$0.03--0.05.

\section{Practical Implications}
\label{sec:implications}

This work shows a useful trade-off for layout-aware discriminative models: they
provide OCR-grounded predictions and operate with 125M parameters on a single
GPU, but the present linker remains substantially below the KVP10k baseline.
A decode-time recovery of unlinked spans extends coverage to the unkeyed and
unvalued categories at markedly lower accuracy (All F1 $0.199\!\to\!0.237$)
without retraining, and the modular pipeline provides a foundation for closing
the remaining gap.

\section{Closing Remarks}
\label{sec:closing}

This thesis demonstrates that effective \kvp{} extraction requires both
layout-aware encoding and appropriate linking granularity. While entity
classification is largely solved by LayoutLMv3 (F1~$= 0.827$), the linking
problem demands careful attention to data pipeline correctness, class
imbalance, and task structure. The progression from V1 through V4 illustrates
that architectural improvements --- such as span-level linking --- can only
realise their potential when the supervision signal is correct. The final
V4 model achieves regular-KVP macro F1~$=0.334$ (text) / $0.309$
(text+location), below the published $0.659/0.611$. The result establishes
feasibility for a compact 125M-parameter discriminative model while identifying
relation modeling and missing singleton-category decoding as remaining work.
